% === IV. SIMULATION ARCHITECTURE =============================================
\section{Simulation Architecture and Data}
\label{sec:platform}

The pipeline runs against a Formula Student vehicle model in
CARLA~\cite{carla2017} through a C\texttt{++} ROS~2 bridge that drives the
vehicle and republishes its state. Two properties of that architecture make
the identifiability claims checkable. CARLA's PhysX solver computes each
wheel's slip angle, normal load and lateral force internally, and the bridge
republishes them on the same clock as the odometry the identifier consumes, so
an identified tire model is scored against the forces the plant generated
rather than against a second model, and that channel never enters the
identification loop. The server also runs synchronously at a fixed
\SI{0.033}{\second} step, with every sample stamped from the simulation clock,
so the one-step residual target is well defined.

The plant is the driverless Formula Student--class vehicle on the
\texttt{silverstone} map (CARLA 0.9.16), whose lateral forces come from the
PhysX tire model~\cite{physxvehicle} under a configured friction coefficient
and lateral stiffness, not from a Magic Formula, so the simulator has no
``true'' Pacejka coefficients to recover and the forces PhysX computes at each
wheel are the ground truth instead. At zero longitudinal slip and zero camber
that model collapses to
\begin{equation}
F_{y}(\alpha)=\mathrm{sgn}(\alpha)\,\mu F_{z}\,S_{1}(K),\qquad
K=\frac{C_{\mathrm{lat}}|\tan\alpha|}{\mu F_{z}},
\label{eq:physx}
\end{equation}
with $C_{\mathrm{lat}}=F_{z}^{\mathrm{rest}}k_{y}S_{1}(3\bar F_{z}/k_{x})$,
$\bar F_{z}$ the normalised load and $S_{1}(K)=\min(1,K-K^{2}/3+K^{3}/27)$ the
engine's smoothing function. It saturates at $\mu F_{z}$, first reached at $K=3$, and never falls off past
that peak, so a Magic Formula can match it up to the peak and no further. Its
scale is set by $\mu$ and $C_{\mathrm{lat}}$, both published by the simulator,
so the plant's curve is read off rather than fitted: $\mu=\num{1.05}$ effective
(a configured wheel friction of \num{1.5} times the road surface's \num{0.70},
with a realised peak axle friction of \numrange{1.00}{1.05}),
$k_{y}=\SI{17}{\per \radian}$, $k_{x}=\num{2.0}$, giving axle slopes of
\num{20.4} and
\SI{19.0}{\kilo\newton\per\radian} front and rear at $\bar F_{z}=1$ and the
static axle loads $F_{z,f}=\SI{1373}{\newton}$, $F_{z,r}=\SI{1274}{\newton}$,
and a peak first reached at \SI{12.0}{\degree} of slip. On 1804 ground-truth
wheel samples from a live lap, each predicted from its own measured $\alpha$,
$F_{z}$ and $\mu$, \eqref{eq:physx} returns $R^{2}=\num{0.99}$ at
\SI{16.1}{\newton} RMSE front and \SI{14.7}{\newton} rear. The vehicle itself
measures $m=\SI{269.6}{\kilo\gram}$ (total static wheel load
\SI{2644.5}{\newton} over $g$, not the configured body mass of
\SI{240.0}{\kilo\gram}), $L=\SI{1.5334}{\meter}$ with
$l_{f}/l_{r}=\num{0.7381}/\SI{0.7953}{\meter}$, an assumed
$I_{z}=\SI{51.1}{\kilo\gram\meter\squared}$, and \SI{16.0}{\degree} of
road-wheel lock.



A model-free pure-pursuit controller drives the vehicle around the reference
raceline. It needs no tire knowledge beyond the wheelbase and is therefore a
valid bootstrap, as in the baseline~\cite{ontrack2025}. Odometry, the achieved
steering angle and the drive command are sampled into a ring buffer for
\SI{30}{\second}, gated on $v_{x}>\SI{1}{\meter\per\second}$, matching the
baseline's data budget, and re-identification retriggers every
\SI{30}{\second}. The model is discretised at $T_{s}=\SI{0.035}{\second}$, a
setting carried because the identification proved insensitive to it:
identifying at \SI{0.020}{\second} returns the same last accepted coefficient
set to four decimal places in the corrected configuration, and still leaves $D$
on its \num{2.0} bound in the inherited one, so the box-edge signature below is
not an artefact of the step mismatch. Preprocessing follows the baseline: a
zero-phase low-pass filter, then symmetry augmentation by mirroring
$(v_{y},\omega,\delta)\rightarrow(-v_{y},-\omega,-\delta)$ at unchanged
$v_{x}$.
